\documentclass[
  spanish,
  aspectratio=169,
  xcolor={dvipsnames,table}
]{beamer}

\geometry{paperwidth=19.2cm,paperheight=10.8cm}

\usepackage{fontspec}
\setmainfont{Arial}
\setsansfont{Arial}
\usepackage[spanish,es-tabla,es-noshorthands]{babel}
\usepackage{array}
\usepackage{booktabs}
\usepackage{graphicx}
\usepackage{ragged2e}
\usepackage{tabularx}
\usepackage{amsmath,amssymb}
\usepackage{tikz}
\usepackage{hyperref}

\definecolor{uanlblue}{RGB}{0,55,108}
\definecolor{uanlgold}{RGB}{218,165,32}
\definecolor{tableblue}{RGB}{233,240,247}
\definecolor{darkcharcoal}{RGB}{33,37,41}

\setbeamercolor{palette primary}{bg=uanlblue,fg=white}
\setbeamercolor{palette secondary}{bg=uanlgold,fg=uanlblue}
\setbeamercolor{structure}{fg=uanlblue}
\setbeamercolor{title}{fg=white,bg=uanlblue}
\setbeamercolor{frametitle}{fg=uanlblue,bg=white}
\setbeamerfont{frametitle}{series=\bfseries,size=\large}

\setbeamertemplate{navigation symbols}{}
\setbeamertemplate{footline}{
  \leavevmode%
  \hbox{%
  \begin{beamercolorbox}[wd=.333333\paperwidth,ht=2.5ex,dp=1.1ex,left,leftskip=1cm]{palette primary}%
    \fontsize{7}{9}\selectfont UANL \textbar\ Facultad de Agronomía
  \end{beamercolorbox}%
  \begin{beamercolorbox}[wd=.333333\paperwidth,ht=2.5ex,dp=1.1ex,center]{palette primary}%
    \fontsize{7}{9}\selectfont Álgebra Lineal — Producto Integrador de Aprendizaje
  \end{beamercolorbox}%
  \begin{beamercolorbox}[wd=.333333\paperwidth,ht=2.5ex,dp=1.1ex,right,rightskip=1cm]{palette primary}%
    \fontsize{7}{9}\selectfont \insertframenumber{} / \inserttotalframenumber
  \end{beamercolorbox}}%
  \vskip0pt%
}

\begin{document}

% -----------------------------------------------------------------------------
% Diapositiva 1: Portada Institucional
% -----------------------------------------------------------------------------
\begin{frame}[plain]
  \begin{tikzpicture}[remember picture,overlay]
    \fill[uanlblue] (current page.south west) rectangle (current page.north east);
    \fill[uanlgold] ([yshift=1.2cm]current page.south west) rectangle ([yshift=1.45cm]current page.south east);
  \end{tikzpicture}
  \vspace{-0.5cm}
  \begin{center}
    \includegraphics[height=1.8cm]{UANL/assets-uanl/logos-institucionales.pdf}
    \vspace{0.35cm}
    
    {\fontsize{10}{12}\selectfont \textbf{\color{uanlgold}UNIVERSIDAD AUTÓNOMA DE NUEVO LEÓN \textbar\ FACULTAD DE AGRONOMÍA}}\\[0.2cm]
    {\fontsize{14}{17}\selectfont \textbf{\color{white}MODELACIÓN MATRICIAL DE FERTILIZACIÓN BALANCEADA Y SISTEMAS LINEALES EN CITRICULTURA}}\\[0.2cm]
    {\fontsize{10}{12}\selectfont \color{white}Ingeniería Agronómica \textbar\ Modelación Matricial de Fertirriego}\\[0.45cm]

    \begin{tabular}{rl}
      \textbf{\color{uanlgold}Estudiante:} & \color{white}Diego de la Cruz Muñoz (Matrícula: 1977875) \\
      \textbf{\color{uanlgold}Docente Titular:} & \color{white}Miranda Vázquez José de Jesús \\
      \textbf{\color{uanlgold}Programa:} & \color{white}Ingeniero Agrónomo \\
      \textbf{\color{uanlgold}Fecha:} & \color{white}11 de noviembre de 2026 \textbar\ General Escobedo / Marín, N.L.
    \end{tabular}
  \end{center}
\end{frame}

% -----------------------------------------------------------------------------
% Diapositiva 2: Introducción y Justificación Agronómica
% -----------------------------------------------------------------------------
\begin{frame}{1. Introducción y Justificación Agronómica}
  \begin{columns}[T]
    \begin{column}{0.48\textwidth}
      \textbf{\color{uanlblue}Problemática Regional en Citricultura:}
      \begin{itemize}
        \item Región citrícola de Nuevo León (General Terán y Montemorelos): alta relevancia socioeconómica en naranja Valencia (\textit{Citrus sinensis}).
        \item La fertilización representa del $30\%$ al $45\%$ de costos directos de operación.
        \item Las fórmulas empíricas tradicionales provocan sobredosificación salina o deficiencias nutrimentales subclínicas.
      \end{itemize}
    \end{column}
    \begin{column}{0.48\textwidth}
      \textbf{\color{uanlblue}Aportación del Álgebra Lineal:}
      \begin{itemize}
        \item \textbf{Modelación determinista:} Transforma el balance químico acoplado en un sistema lineal $A\mathbf{x} = \mathbf{b}$.
        \item \textbf{Cero margen de error:} Permite obtener la dosis física exacta ($\text{kg/ha}$) sin excesos lixiviables ni déficit.
        \item \textbf{Sustentabilidad y rentabilidad:} Protege los mantos freáticos y optimiza el costo financiero de los insumos.
      \end{itemize}
    \end{column}
  \end{columns}
\end{frame}

% -----------------------------------------------------------------------------
% Diapositiva 3: Recolección de Datos
% -----------------------------------------------------------------------------
\begin{frame}{2. Recolección de Datos y Parámetros del Sistema}
  \begin{columns}[T]
    \begin{column}{0.48\textwidth}
      \textbf{\color{uanlblue}Demanda Neta por Hectárea (Meta: $35\text{ t/ha}$):}
      \begin{itemize}
        \item \textbf{Nitrógeno ($N$):} $b_1 = 175\text{ kg/ha}$
        \item \textbf{Fósforo ($P_2O_5$):} $b_2 = 142\text{ kg/ha}$
        \item \textbf{Potasio ($K_2O$):} $b_3 = 182\text{ kg/ha}$
      \end{itemize}
      \vspace{0.2cm}
      \textbf{\color{uanlblue}Insumos Comerciales Acoplados:}
      \begin{itemize}
        \item $x_1$: NPK 20-10-10 ($20\% N, 10\% P_2O_5, 10\% K_2O$).
        \item $x_2$: DAP 18-46-0 ($18\% N, 46\% P_2O_5, 0\% K_2O$).
        \item $x_3$: Nitrato Potásico 13-0-44 ($13\% N, 0\% P_2O_5, 44\% K_2O$).
      \end{itemize}
    \end{column}
    \begin{column}{0.48\textwidth}
      \textbf{\color{uanlblue}Matriz de Concentración Nutrimental:}
      \vspace{0.15cm}
      \small
      \begin{tabular}{lccc}
        \toprule
        \textbf{Nutriente} & \textbf{NPK ($x_1$)} & \textbf{DAP ($x_2$)} & \textbf{KNO$_3$ ($x_3$)} \\
        \midrule
        $N$ & $0.20$ & $0.18$ & $0.13$ \\
        $P_2O_5$ & $0.10$ & $0.46$ & $0.00$ \\
        $K_2O$ & $0.10$ & $0.00$ & $0.44$ \\
        \bottomrule
      \end{tabular}
      \vspace{0.3cm}
      \begin{block}{Interdependencia Química}
        Ninguna dosis puede ajustarse de forma aislada sin perturbar el balance total de macronutrientes.
      \end{block}
    \end{column}
  \end{columns}
\end{frame}

% -----------------------------------------------------------------------------
% Diapositiva 4: Formulación y Compatibilidad Matricial
% -----------------------------------------------------------------------------
\begin{frame}{3. Formulación del Sistema y Criterio de Rouché-Capelli}
  \begin{block}{Formulación Matricial Canónica: $A\mathbf{x} = \mathbf{b}$}
    $$\begin{pmatrix}
        0.20 & 0.18 & 0.13 \\
        0.10 & 0.46 & 0.00 \\
        0.10 & 0.00 & 0.44
      \end{pmatrix}
      \begin{pmatrix} x_1 \\ x_2 \\ x_3 \end{pmatrix}
      =
      \begin{pmatrix} 175 \\ 142 \\ 182 \end{pmatrix}$$
  \end{block}
  \vspace{0.15cm}
  \begin{columns}[T]
    \begin{column}{0.48\textwidth}
      \textbf{\color{uanlblue}Cálculo del Determinante:}
      $$\det(A) = 0.20(0.46 \times 0.44) - 0.18(0.10 \times 0.44) + 0.13(-0.046)$$
      $$\det(A) = 0.04048 - 0.00792 - 0.00598 = \mathbf{0.02658} \neq 0$$
    \end{column}
    \begin{column}{0.48\textwidth}
      \textbf{\color{uanlblue}Diagnóstico Estructural:}
      \begin{itemize}
        \item $\det(A) \neq 0 \implies A$ es regular e invertible.
        \item $\operatorname{rango}(A) = \operatorname{rango}[A|\mathbf{b}] = n = 3$.
        \item \textbf{Conclusión:} El sistema es \textit{compatible determinado}; existe una única solución física.
      \end{itemize}
    \end{column}
  \end{columns}
\end{frame}

% -----------------------------------------------------------------------------
% Diapositiva 5: Procedimiento 1: Gauss-Jordan
% -----------------------------------------------------------------------------
\begin{frame}{4. Resolución por Eliminación de Gauss-Jordan}
  \begin{block}{Pivoteo sobre la Matriz Aumentada $[A|\mathbf{b}]$}
    $$[A|\mathbf{b}]^{(0)} = \left(\begin{array}{ccc|c}
      0.20 & 0.18 & 0.13 & 175 \\
      0.10 & 0.46 & 0.00 & 142 \\
      0.10 & 0.00 & 0.44 & 182
    \end{array}\right)
    \xrightarrow{R_1 \leftarrow 5R_1}
    \left(\begin{array}{ccc|c}
      1.00 & 0.90 & 0.65 & 875 \\
      0.10 & 0.46 & 0.00 & 142 \\
      0.10 & 0.00 & 0.44 & 182
    \end{array}\right)$$
  \end{block}
  \begin{itemize}
    \item \textbf{Anulación de columna 1:} $R_2 \leftarrow R_2 - 0.1R_1$, $R_3 \leftarrow R_3 - 0.1R_1 \implies [0, 0.37, -0.065 \mid 54.50]$ y $[0, -0.09, 0.375 \mid 94.50]$.
    \item \textbf{Normalización de pivote 2 y eliminación:} $R_2 \leftarrow \frac{100}{37}R_2$, $R_1 \leftarrow R_1 - 0.9R_2$, $R_3 \leftarrow R_3 + 0.09R_2$.
    \item \textbf{Pivote 3 y reducción regresiva:} $a_{33}' = \frac{1329}{3700}$, $b_3' = \frac{3987}{37} \implies x_3 = \frac{3987}{37} \times \frac{3700}{1329} = \mathbf{300\text{ kg/ha}}$.
    \item Sustituyendo hacia arriba: $x_2 = \mathbf{200\text{ kg/ha}}$ y $x_1 = \mathbf{500\text{ kg/ha}}$.
  \end{itemize}
\end{frame}

% -----------------------------------------------------------------------------
% Diapositiva 6: Procedimiento 2: Regla de Cramer
% -----------------------------------------------------------------------------
\begin{frame}{5. Corroboración por Regla de Cramer}
  \begin{columns}[T]
    \begin{column}{0.48\textwidth}
      \textbf{\color{uanlblue}Determinantes Modificados:}
      \begin{itemize}
        \item $\det(A) = 0.02658$
        \item $\det(A_1) = \det\begin{pmatrix} 175 & 0.18 & 0.13 \\ 142 & 0.46 & 0.00 \\ 182 & 0.00 & 0.44 \end{pmatrix} = \mathbf{13.29}$
        \item $\det(A_2) = \det\begin{pmatrix} 0.20 & 175 & 0.13 \\ 0.10 & 142 & 0.00 \\ 0.10 & 182 & 0.44 \end{pmatrix} = \mathbf{5.316}$
        \item $\det(A_3) = \det\begin{pmatrix} 0.20 & 0.18 & 175 \\ 0.10 & 0.46 & 142 \\ 0.10 & 0.00 & 182 \end{pmatrix} = \mathbf{7.974}$
      \end{itemize}
    \end{column}
    \begin{column}{0.48\textwidth}
      \textbf{\color{uanlblue}Dosis Exactas Obtenidas:}
      \begin{align*}
        x_1 &= \frac{\det(A_1)}{\det(A)} = \frac{13.29}{0.02658} = \mathbf{500\text{ kg/ha}} \\[0.2cm]
        x_2 &= \frac{\det(A_2)}{\det(A)} = \frac{5.316}{0.02658} = \mathbf{200\text{ kg/ha}} \\[0.2cm]
        x_3 &= \frac{\det(A_3)}{\det(A)} = \frac{7.974}{0.02658} = \mathbf{300\text{ kg/ha}}
      \end{align*}
      \begin{block}{Coincidencia Total}
        Cramer valida con exactitud analítica simbólica la solución de Gauss-Jordan.
      \end{block}
    \end{column}
  \end{columns}
\end{frame}

% -----------------------------------------------------------------------------
% Diapositiva 7: Procedimiento 3: Matriz Inversa y Residuo
% -----------------------------------------------------------------------------
\begin{frame}{6. Matriz Inversa y Verificación de Residuo}
  \begin{columns}[T]
    \begin{column}{0.50\textwidth}
      \textbf{\color{uanlblue}Inversión Matricial: $\mathbf{x} = A^{-1}\mathbf{b}$}
      $$A^{-1} = \frac{1}{0.02658} \begin{pmatrix}
        0.2024 & -0.0792 & -0.0598 \\
        -0.0440 & 0.0750 & 0.0130 \\
        -0.0460 & 0.0180 & 0.0740
      \end{pmatrix}$$
      $$\mathbf{x} = A^{-1} \begin{pmatrix} 175 \\ 142 \\ 182 \end{pmatrix} = \begin{pmatrix} \mathbf{500} \\ \mathbf{200} \\ \mathbf{300} \end{pmatrix}$$
    \end{column}
    \begin{column}{0.46\textwidth}
      \textbf{\color{uanlblue}Comprobación del Residuo $\mathbf{r} = A\mathbf{x} - \mathbf{b}$:}
      \begin{align*}
        r_1 &= 0.20(500) + 0.18(200) + 0.13(300) - 175 \\
            &= (100 + 36 + 39) - 175 = \mathbf{0} \\[0.15cm]
        r_2 &= 0.10(500) + 0.46(200) + 0 - 142 \\
            &= (50 + 92) - 142 = \mathbf{0} \\[0.15cm]
        r_3 &= 0.10(500) + 0 + 0.44(300) - 182 \\
            &= (50 + 132) - 182 = \mathbf{0}
      \end{align*}
      \textbf{Residuo nulo:} $\|\mathbf{r}\|_\infty = 0$.
    \end{column}
  \end{columns}
\end{frame}

% -----------------------------------------------------------------------------
% Diapositiva 8: Resultados y Dosificación Final
% -----------------------------------------------------------------------------
\begin{frame}{7. Resultados: Dosificación Agronómica de Campo}
  \begin{table}[htbp]
    \centering
    \small
    \begin{tabularx}{0.92\textwidth}{lcccc}
      \toprule
      \textbf{Fuente comercial} & \textbf{Dosis (kg/ha)} & \textbf{Aporte $N$} & \textbf{Aporte $P_2O_5$} & \textbf{Aporte $K_2O$} \\
      \midrule
      Mezcla NPK 20-10-10 & $500$ ($10$ sacos) & $100$ & $50$ & $50$ \\
      Fosfato Diamónico (DAP) & $200$ ($4$ sacos) & $36$ & $92$ & $0$ \\
      Nitrato de Potasio & $300$ ($6$ sacos) & $39$ & $0$ & $132$ \\
      \midrule
      \textbf{Total Suministrado} & $\mathbf{1000\text{ kg/ha}}$ & $\mathbf{175\text{ kg/ha}}$ & $\mathbf{142\text{ kg/ha}}$ & $\mathbf{182\text{ kg/ha}}$ \\
      \textbf{Meta Nutrimental} & --- & $\mathbf{175\text{ kg/ha}}$ & $\mathbf{142\text{ kg/ha}}$ & $\mathbf{182\text{ kg/ha}}$ \\
      \bottomrule
    \end{tabularx}
  \end{table}
  \begin{center}
    \textbf{\color{uanlblue}Balance exacto al $100\%$ sin desperdicio ni toxicidad edáfica.}
  \end{center}
\end{frame}

% -----------------------------------------------------------------------------
% Diapositiva 9: Análisis de Estabilidad y Fertirriego
% -----------------------------------------------------------------------------
\begin{frame}{8. Estabilidad Numérica y Plan de Fertirriego}
  \begin{columns}[T]
    \begin{column}{0.48\textwidth}
      \textbf{\color{uanlblue}Número de Condición: $\kappa(A) \approx 5.86$}
      \begin{itemize}
        \item $\kappa(A) \ll 100 \implies$ Matriz \textit{bien condicionada}.
        \item Inmune a variaciones menores en la ley comercial de los fertilizantes ($\pm 2\%$).
        \item Estabilidad garantizada durante la aplicación práctica en campo.
      \end{itemize}
    \end{column}
    \begin{column}{0.48\textwidth}
      \textbf{\color{uanlblue}Plan de Distribución Fenológica (36 semanas):}
      \begin{itemize}
        \item \textbf{Fase 1: Brotación floral (15\%):} $150\text{ kg/ha}$ (prioridad $P_2O_5$).
        \item \textbf{Fase 2: Cuajado y división (50\%):} $500\text{ kg/ha}$ (equilibrio $N$-$P$).
        \item \textbf{Fase 3: Llenado de fruto (35\%):} $350\text{ kg/ha}$ (prioridad $K_2O$).
      \end{itemize}
    \end{column}
  \end{columns}
\end{frame}

% -----------------------------------------------------------------------------
% Diapositiva 10: Conclusiones
% -----------------------------------------------------------------------------
\begin{frame}{9. Conclusiones y Transferencia Profesional}
  \begin{enumerate}
    \item \textbf{Rigor analítico:} El álgebra lineal elimina el error empírico en la fertilización de cítricos en Nuevo León mediante sistemas acoplados $3 \times 3$.
    \item \textbf{Convergencia de métodos:} Gauss-Jordan, Cramer y la Matriz Inversa demostraron consistencia absoluta con residuo nulo ($\mathbf{r} = \mathbf{0}$).
    \item \textbf{Sostenibilidad edáfica:} La dosificación de $500\text{ kg}$ de NPK, $200\text{ kg}$ de DAP y $300\text{ kg}$ de KNO$_3$ optimiza costos y previene la salinización secundaria.
    \item \textbf{Competencia profesional:} La modelación matemática capacita al agrónomo de la UANL para liderar decisiones predictivas basadas en datos rigurosos.
  \end{enumerate}
\end{frame}

% -----------------------------------------------------------------------------
% Diapositiva 11: Referencias Bibliográficas
% -----------------------------------------------------------------------------
\begin{frame}{Referencias Bibliográficas}
  \fontsize{8}{10}\selectfont
  \begin{itemize}
    \item Anton, H., \& Rorres, C. (2014). \textit{Elementary Linear Algebra: Applications Version} (11th ed.). John Wiley \& Sons.
    \item Grossman, S. I., \& Flores Godoy, J. J. (2019). \textit{Álgebra lineal} (8va ed.). McGraw-Hill Education.
    \item Lay, D. C., Lay, S. R., \& McDonald, J. J. (2016). \textit{Linear Algebra and Its Applications} (5th ed.). Pearson.
    \item Strang, G. (2016). \textit{Introduction to Linear Algebra} (5th ed.). Wellesley-Cambridge Press.
    \item University of Minnesota Extension. (2024). \textit{Nutrient Management and Commercial Fertilizer Recommendations}. UMN Extension Publications.
  \end{itemize}
\end{frame}

\end{document}
